% !TEX TS-program = lualatexmk
% !TEX parameter =  --shell-escape
\documentclass[benediction_assumption_la_en.tex]{subfiles}

\ifcsname preamble@file\endcsname
  \setcounter{page}{\getpagerefnumber{M-benediction_assumption_la_en_other}}
\fi

\begin{document}

\chapter*{OTHER RITES.}
\addcontentsline{toc}{chapter}{Other Rites}\phantomsection

\smalltitle{Holy Communion distributed outside of Mass on First Fridays or on other occasions.}
\addcontentsline{toc}{section}{Holy Communion distributed outside of Mass} \phantomsection

\rubrique{\begin{enpars}This rite is said entirely in Latin: the translation is for the convenience of the faithful.\end{enpars}}

\rubrique{
\begin{enpars}
At least two candles are lit. The priest, wearing a surplice with a white stole or one in the color of the day (on All Souls, white or violet), or the Mass vestments when distributing communion just before or after Mass, approaches the altar with an acolyte. The priest spreads the corporal contained in the burse while the server, joined by the faithful in a reverent manner, begins the \normaltext{Confitéor,} bowing from the waist while kneeling at the epistle corner of the altar.\end{enpars}}

\rubrique{The rite is celebrated identically when a deacon distributes communion, except that he wears a diaconal stole.}

\textes{confiteor_la_en}

\rubrique{\begin{enpars}The priest, having removed the ciborium lid, genuflects and turns a little towards the people, and says:\end{enpars}}

\textes{misereatur_vestri_la_en}

\rubrique{\begin{enpars}The priest turns back to the altar, genuflects, and takes the ciborium in his hands. With a host held between his thumb and index and over the ciborium, he then turns towards the people, saying:\end{enpars}}

\texteparacol{\lettrine{E}{c}ce Agnus Dei, ecce qui tollit peccáta mundi.}{english}{\noindent Behold the Lamb of God, behold him who taketh away the sins of the world.}

\rubrique{Then the following is said three times by the priest together with the faithful.}

\texteparacol{\lettrine{D}{o}mine, non sum dignus, ut intres
sub tectum meum; sed tantum dic
verbo, et sanábitur ánima mea.}{english}{\noindent Lord, I am not worthy that thou shouldst enter under my roof, but only say the word, and my soul shall be healed.}

\rubrique{\begin{enpars}Then the priest goes to the communion rail to distribute Holy Communion, beginning at the epistle end, making the sign of the cross with each particle and
saying to each communicant:\end{enpars}}

\texteparacol{\lettrine{C}{o}rpus \cc{} Dómini nostri Jesu Christi
custódiat ánimam tuam in vitam ætérnam. Amen.}{english}{\noindent May the Body of our Lord Jesus Christ preserve thy soul unto life everlasting. Amen.}

\rubrique{Where \normaltext{Allelúia} is to be said in Paschal Time, it is also said during the octave of Corpus Christi.}

\texteparacol{\lettrine{O}{} Sacrum convívium, in quo Christus
súmitur, recólitur memória passiónis
eius, mens implétur grátia, et futúræ
glóriæ nobis pignus datur. \tpalleluia.

\vv Panem de cælo præstitísti eis. \tpalleluia

\rr Omne dele\-ctaméntum in se habéntem. \tpalleluia

\vv Dómine, exaudi oratiónem meam.

\rr Et clamor meus ad te véniat.

\vv Dóminus vobíscum.

\rr Et cum spíritu tuo.

Orémus.
  
  \lettrine{D}{e}us, qui nobis sub Sacraménto mirábili passiónis tuæ memóriam reliquísti: tríbue, quǽsumus, ita nos córporis, et sánguinis tui sacra mystéria venerári; ut redem\-ptiónis tuæ fru\-ctum in nobis júgiter sentiámus. Qui vivis et regnas cum Deo patre in sǽcula sæculórum in unitáte Spíritus San\-cti Deus, per ómnia sǽcula sæculórum.
  
  \rr Amen.}
{english}{\noindent O Sacred Banquet, in which Christ is received, the memory of his Passion is recalled, the soul is filled with grace, and the pledge of future glory is given unto us. \textit{P.T.} Alleluia.

\noindent \vv Thou hast given them bread from heaven. (\textit{P.T.} alleluia.)

\noindent \rr Having all sweetness within it. (\textit{P.T.} alleluia.)

\noindent \vv O Lord, hear my prayer.

\noindent \rr And let my prayer come unto thee.

\noindent \vv The Lord be with you.

\noindent \rr And with thy spirit.

\noindent Let us pray. O God, who under a wonderful Sacrament hast left us a memorial of thy Passion; grant us, we beseech thee, so to reverence the sacred mysteries of thy Body and Blood, that we may ever feel within ourselves the fruit of thy Redemption. Who livest and reignest with God the Father in the unity of the Holy Ghost, world without end. Amen.}

\rubrique{In Paschal time, the following collect is said instead.}

 \texteparacol{Orémus.
 
\lettrine{S}{p}íritum nobis, Dómine, tuæ caritátis infúnde: ut, quos sacraméntis paschálibus satiásti, tua fácias pietáte concórdes.
Per Dóminum nostrum Jesum Christum, Fílium tuum: qui tecum vivit et regnat in unitáte ejúsdem Spíritus San\-cti Deus, per ómnia sǽcula sæculórum.

\rr Amen.}{english}{\noindent Let us pray. Pour forth, we beseech thee o Lord, the spirit of love into our hearts, that thou by thy mercy 
might make to be of one mind those whom thou hast satisfied with the paschal sacraments. Through Jesus Christ, thy Son our Lord, who liveth and reigneth with thee, in the unity of the same Holy Ghost, God, world without end. Amen.}

%{}\footnote{Atque per octavam Corporis Christi. \textit{As well as throughout the octave of Corpus Christi.}

\rubrique{While saying the prayers, \normaltext{O sacrum
convívium,} etc., the priest purifies first the thumb and forefinger, then the communion plate over the
ciborium. With the thumb and forefinger still
joined, he covers the ciborium and then purifies his fingers in
the ablution cup and dry them with the purificator.}

\rubrique{The priest opens the tabernacle, replaces the
ciborium, genuflects, closes and locks the tabernacle door.}

\rubrique{Then, he brings his hands apart and then together again, turns to the people, he blesses them, keeping his left hand on his chest.}

\texteparacol{\lettrine{B}{e}nedí\-ctio Dei omnipoténtis, Patris,
et Fílii, \cc{} et Spíritus San\-cti, descéndat super vos, et máneat semper.

\rr Amen.}{english}{\noindent May the blessing of almighty God, Father, Son, and Holy Ghost, descend upon you and remain with you forever.}

\rubrique{Then \normaltext{Cor Jesu sacratíssimum, \pageref{M-cor_jesu},} or another suitable chant follows.}

\rubrique{The rite is celebrated identically by a bishop except for the pontifical blesssing.}

%%noindent preference poses problem with lettrine here
%%sligh misalignment and syllable break seems preferable
\texteparacol{\lettrine{S}{i}t nomen Dómini benedí\-ctum.\\
\rr Ex hoc nunc et usque in sǽculum.

\vv Adjutórium \cc{} nostrum in nómine Dómini.

\rr Qui fecit cælum et terram.

\lettrine{B}{e}nedí\-ctio Dei omnipoténtis, Patris, \cc{}
et Fílii, \cc{} et Spíritus \cc{} San\-cti, descéndat super vos, et máneat semper.

\rr Amen.}{english}
{\noindent \vv Blessed be the name of the Lord.

\noindent \rr Now and forever.

\noindent \vv Our help is in the name of the Lord.

\noindent \rr Who made heaven and earth.

\noindent \vv May the blessing of almighty God, Father, Son, and Holy Ghost, descend upon you and remain with you forever.

\noindent \rr Amen.}

\rubrique{If black vestments are worn when distributing communion just before or after Mass, the blessing is omitted, as is \normaltext{Allelúia} in Paschal Time.}

\newpage

\smalltitle{The Major Litanies on the feast of Saint Mark and the Minor Litanies before the Ascension}
\addcontentsline{toc}{section}{For processions.} \phantomsection

{\centering{and in other processions.\footnote{The special invocations of the litany and the prayers for other occasions are found in the ritual.}\par}}

\rubrique{Ante Processionem, cantatur stando. Before the procession, the following is sung, standing.}

\gscore[Ant.]{2.}{an_exsurge_domine_adjuva_en_solesmes}

\rubrique{Duo cantores ante altare majus genuflexi Litanias incipiunt, et omnes versus duplicantur, nisi Processio fuerit impedita. Surgunt omnes et ordinatim procendunt, ad
\normaltext{Sancta Dei}.
}

\rubrique{Two cantors kneeling before the main altar begin the Litany. All verses are doubled, except if the procession is impeded. All rise and the procession departs at \normaltext{Sancta Dei.}}

\rubrique{\normaltext{Litaniæ Sanctorum, \pageref{litaniae_sanctorum},} ps.\@ 69 et preces.}

\rubrique{\normaltext{Litany of the Saints, \pageref{litaniae_sanctorum}.} with ps.\@ 69 and the prayers.}

\newpage

\smalltitle{Forty Hours'\thinspace Devotion.}

\addcontentsline{toc}{section}{Forty Hours'\thinspace Devotion}\phantomsection
%. After the Mass on the day of exposition.
\texteparacol{In die expositionis, missa finita, Celebrans ad scammum deponit planetam et
manipulum et assumit pluviale. S. Ministri deponunt manipula.
 Celebrans ad scammum incensum in duo thuribula sine benedictione injicit.
 
Celebrans genuflexus in infimo gradu Altaris SS.\@ Sacramentum incensat. Deinde assumit velum humerale.
 
Fit processio. Intonatur \textit{Pange lingua,} ut infra.
 
Peracta processione, Diaconus SS. Sacramentum in throno collocat, interim Celebrans deponit velum humerale, et deinde intonatur
\textit{Tantum ergo Sacramentum.} Ad \textit{Genitori} incensatur SS.\@ Sacramentum
more solito.

Completo cantu hymni \textit{Tantum ergo,} non additur \textit{Panem de cælo,}
etc., sed statim Cantores incipiunt Litanias Sanctorum.}
{english}
{After Mass on the day of exposition, the celebrant removes the chasuble and maniple
at the sedilia and puts on the cope. The sacred ministers remove their maniples. The celebrant imposes incense without blessing it in the two thuribles.

The celebrant, kneeling on the lowest grade of the altar steps, incenses the Most Blessed Sacrament, then he puts on the humeral veil.

The procession takes place. \textit{Pangue lingua} is intoned, as below.

At the end of the procession, the deacon places the Most Blessed Sacrament on the throne, while the celebrant removes the humeral veil. Then \textit{Tantum ergo Sacramentum} is intoned. At \textit{Genitori,} the Most Blessed Sacrament is incensed in the usual way.

After the hymn \textit{Panem de cælo} is not said, but the cantors immediately begin the Litany of the Saints.}

\label{hy_pange_lingua_corpus_christi_solesmes} \phantomsection
\gscore[]{3.}{hy_pange_lingua_corpus_christi_solesmes}
\textes{pange_lingua_gloriosi_en}

\myrule
\rubrique{In die repositionis tantum. On the day of reposition only.}
\texteparacol{\vv Panem de cælo præstitísti eis. \tpalleluia{}\footnote{Atque per octavam Corporis Christi. \textit{As well as throughout the octave of Corpus Christi.}}

\rr Omne dele\-ctaméntum in se habéntem. \tpalleluia{}}
{english}
{\noindent \vv Thou hast given them bread from heaven. (\textit{P.T.} Alleluia.)

\noindent \rr Having within it all sweetness. (\textit{P.T.} Alleluia.)}

\myrule

\textes{xl_horarum_orationes}

\texteparacol{In die repositionis, post Missam Celebrans, ad scammum, deponit casulam et
manipulum et assumit pluviale. S. Ministri manipula deponunt, et omnes ante Altare redeunt.

Cantantur Litaniæ Sanctorum et preces usque ad ℣. \textit{Dómine exáudi oratiónem meam}
inclusive.

Surgit Celebrans et incensum in duo thuribula sine benedictione injicit, SS.\@ Sacramentum incensat, et assumit velum humerale.

Fit processio. Intonatur \textit{Pange lingua,} \pageref{hy_pange_lingua_corpus_christi_solesmes}.

Peracta processione, et SS.\@ Sacramento super Altare collocato, Celebrans deponit velum humerale.

Intonatur
\textit{Tantum ergo.} Ad \textit{Genitori,} Celebrans incemsum injicit et incensatur SS.\@ Sacramentum.

Completo cantu hymni \textit{Tantum ergo,} intonatur \textit{Panem de cælo,} etc.

Surgit Celebrans et, quin præmittat \textit{Dóminus vobíscum,} cantat omnes orationes præscri\-ptas.

Quibus absolutis rursus genuflectit et cantat \textit{Dómine exáudi} etc. Deinde cantores cantant \textit{Exáudiat nos,} etc., et Celebrans submissa voce addit \textit{Et fidélium,} etc.

Celebrans assumit velum humerale, et Benedictionem impertit more solito.}
{english}
{After Mass on the day of reposition, the celebrant removes the chasuble and maniple
at the sedilia and puts on the cope. The sacred ministers remove their maniples.

The Litany of the Saints is sung with its prayers until ℣.
\textit{O Lord, hear my prayer} inclusively.

The celebrant rises and imposes incense without blessing it in the two thuribles, and then he puts on the humeral veil.

The procession takes place. \textit{Pangue lingua} is intoned, \pageref{hy_pange_lingua_corpus_christi_solesmes}.

At the end of the procession, the Most Blessed Sacrament having been replaced on the throne, the celebrant removes the humeral veil. Then \textit{Tantum ergo Sacramentum} is intoned. At \textit{Genitori,} the celebrant imposes incense and Most Blessed Sacrament is incensed in the usual way.

Having finished the hymn \textit{Tantum ergo,} \textit{Panem de cælo,} etc. are sung.

The celebrant rises, and without saying \textit{Dóminus vobíscum,} sings the prescribed prayers.

Having finished these, he kneels again and sings \textit{Dómine exáudi} etc. Then the cantors sing
\textit{Exáudiat nos,} etc., and the celebrant adds in a lower voice. \textit{Et fidélium,} etc.

The celebrant puts on the humeral veil again, and Benediction is given in the usual way.}

%\newpage
%
%\newpage

%%litanies align with odd page as first one

\cleardoublepage

\label{litaniae_sanctorum} \phantomsection
\smalltitle{Litany of the Saints.}

\smallscore{litaniae_sanctorum_1_kyrie_eleison_solesmes}

\gresetlastline{justified}
\smallscore{litaniae_sanctorum_2_pater_solesmes}%
\begin{nstabbing}
\>Fili \>Re-\>dém-\>ptor mundi \>\>\textbf{De}-\>us, \>mi-\>se-\>\textit{ré}-\>\textit{re} \>no-\>bis.\\
%\>Fili \>Re-\>démp-\>tor \>mundi \>\textbf{De}-\>us, \>mi-\>se-\>\textit{ré}-\>\textit{re} \>no-\>bis.\\
\>Spí-\>ri-\>tus \>Sancte \>\>\textbf{De}-\>us, \>mi-\>se-\>\textit{ré}-\>\textit{re} \>no-\>bis.\\
\>Sancta \>\>Trínitas \>\>unus \>\textbf{De}-\>us, \>mi-\>se-\>\textit{ré}-\>\textit{re} \>no-\>bis.\\
\end{nstabbing}

%%can we improve alignment of dactyls?
%%2nd syllables are all misaligned
%%only vir- is aligned
%%last are also wrong, but it's subtly in between e and l for Angels
%%however it resembles example http://gregorio-project.github.io/tips/litany.html
%%so for now we will call it a day
\smallscore{litaniae_sanctorum_3_sancta_maria_solesmes}
\begin{nstabbing}
\>Sancta Dé-\>\>i \>\textbf{Gé}-\>ni- \>trix, \>ó-\>\textit{ra} \>\textit{pro} \>nó-\>bis.\\
\>Sancta Vír-\>\>go \>\textbf{vír-}\>gi-\>num, \>ó-\>\textit{ra} \>\textit{pro} \>nó-\>bis.\\
\>Sancte \>\>\>\textbf{Mí-}\>cha-\>el, \>ó-\>\textit{ra} \>\textit{pro} \>nó-\>bis.\\
\>Sancte \>\>\>\textbf{Gá-}\>bri-\>el, \>ó-\>\textit{ra} \>\textit{pro} \>nó-\>bis.\\
\>Sancte \>\>\>\textbf{Rá-}\>pha-\>el, \>ó-\>\textit{ra} \>\textit{pro} \>nó-\>bis.\\
\end{nstabbing}

\smallscore{litaniae_sanctorum_4_omnes_solesmes}
%\noindent
\begin{nstabbing}
\>Omnes \quad sancti \quad beatórum Spirítuum \>\>\>\>\>\textbf{ór}di-nes, \>\>o-\>rá-\>\textit{te} \>\textit{pro} \>nó-\>bis.\\
\>Sancte \>Joánnes \>\>Ba\>\textbf{ptí}-\>\>sta, \>\>o-\>\textit{ra} \>\textit{pro} \>nó-\>bis.\\
\>Sancte \>\>\>\>\>\textbf{Jó-}\>seph, \>\>o-\>\textit{ra} \>\textit{pro} \>nó-\>bis.\\
\>Omnes \>sancti \>Patriárchæ et \quad \>Pro-\>\>\textbf{phé}-\>tæ, \>o-\>rá-\>\textit{te} \>\textit{pro} \>nó-\>bis.\\
\end{nstabbing}
\gresetlastline{ragged}
\smallscore{litaniae_sanctorum_5_sancte_petre_solesmes}
\begin{multicols}{2}
\noindent
San\-cte \textbf{Pau}le,\hfill o\textit{ra}.\\
San\-cte An\textbf{dré}a,\hfill o\textit{ra}.\\
San\-cte Ja\textbf{có}be,\hfill o\textit{ra}.\\
San\-cte Jo\textbf{án}nes,\hfill o\textit{ra}.\\
San\-cte \textbf{Tho}ma,\hfill o\textit{ra}.\\
San\-cte Ja\textbf{có}be,\hfill o\textit{ra}.\\
San\-cte Phi\textbf{líp}pe,\hfill o\textit{ra}.\\
San\-cte Bartholo\textbf{mǽ}e,\hfill o\textit{ra}.\\
San\-cte Mat\textbf{thǽ}e,\hfill o\textit{ra}.\\
San\-cte \textbf{Si}mon,\hfill o\textit{ra}.\\
San\-cte Thad\textbf{dǽ}e,\hfill o\textit{ra}.\\
San\-cte Mat\textbf{thí}a,\hfill o\textit{ra}.\\
San\-cte \textbf{Bár}naba,\hfill o\textit{ra}.\\
San\-cte \textbf{Lu}ca,\hfill o\textit{ra}.\\
San\-cte \textbf{Mar}ce,\hfill o\textit{ra}.\\
Omnes san\-cti Apóstoli et \\ \emspace Evange\textbf{lí}stæ,\hfill orá\textit{te}.\\
Omnes san\-cti Discípuli \textbf{Dó}mini,\hfill orá\textit{te}.\\
Omnes san\-cti Inno\textbf{cén}tes,\hfill orá\textit{te}.\\
San\-cte \textbf{Sté}phane,\hfill o\textit{ra}.\\
San\-cte Lau\textbf{rén}ti,\hfill o\textit{ra}.\\
San\-cte Vin\textbf{cén}ti,\hfill o\textit{ra}.\\
San\-cti Fabiáne et Sebasti\textbf{án}e,\hfill orá\textit{te}.\\
San\-cti Joánnes et \textbf{Pau}le,\hfill orá\textit{te}.\\
San\-cti Cosma et Dami\textbf{án}e,\hfill orá\textit{te}.\\
San\-cti Gervási et Pro\textbf{tá}si,\hfill orá\textit{te}.\\
Omnes san\-cti \textbf{Már}tyres,\hfill orá\textit{te}.\\
San\-cte Sil\textbf{vé}ster,\hfill o\textit{ra}.\\
San\-cte Gre\textbf{gó}ri,\hfill o\textit{ra}.\\
San\-cte Am\textbf{bró}si,\hfill o\textit{ra}.\\
San\-cte Augus\textbf{tí}ne,\hfill o\textit{ra}.\\
San\-cte Hie\textbf{ró}nyme,\hfill o\textit{ra}.\\
San\-cte Mar\textbf{tí}ne,\hfill o\textit{ra}.\\
San\-cte Nico\textbf{lá}e,\hfill o\textit{ra}.\\
Omnes san\-cti Pontífices et \\ \emspace Confes\textbf{só}res,\hfill orá\textit{te}.\\
Omnes san\-cti Do\textbf{ctó}res,\hfill orá\textit{te}.\\
San\-cte An\textbf{tó}ni,\hfill o\textit{ra}.\\
San\-cte Bene\textbf{dí}cte,\hfill o\textit{ra}.\\
San\-cte Ber\textbf{nár}de,\hfill o\textit{ra}.\\
San\-cte Do\textbf{mí}nice,\hfill o\textit{ra}.\\
San\-cte Fran\textbf{cí}sce,\hfill o\textit{ra}.\\
Omnes san\-cti Sacerdótes et \\  \emspace Le\textbf{ví}tæ,\hfill orá\textit{te}.\\
Omnes san\-cti Mónachi et \\ \emspace Ere\textbf{mí}tæ,\hfill orá\textit{te}.\\
San\-cta María Magda\textbf{lé}na,\hfill o\textit{ra}.\\
San\-cta \textbf{A}gatha,\hfill o\textit{ra}.\\
San\-cta \textbf{Lú}cia,\hfill o\textit{ra}.\\
San\-cta \textbf{A}gnes,\hfill o\textit{ra}.\\
San\-cta Cæ\textbf{cí}lia,\hfill o\textit{ra}.\\
San\-cta Catha\textbf{rí}na,\hfill o\textit{ra}.\\
San\-cta Anas\textbf{tá}sia,\hfill o\textit{ra}.\\
Omnes san\-ctæ Vírgines \\ \emspace et \textbf{Ví}duæ,
\hfill orá\textit{te}.\\
Omnes San\-cti et San\-ctæ \textbf{De}i,\\ 
\zeroptspace \hfill intercédi\textit{te pro} nobis.
\end {multicols}

%%less than ideal formatting shared with gregobase
\smallscore{litaniae_sanctorum_6_propitius_esto_solesmes}%

\noindent
Propí\textit{tius} \textbf{es}to,\hfill exáudi nos Dómine.\\
Ab \textit{omni} \textbf{ma}lo,\hfill líbera nos Dómine.\\
Ab \textit{omni pec}\textbf{cá}to,\hfill líbera nos Dómine.\\
Ab \textit{ira} \textbf{tu}a,\hfill líbera nos Dómine.\\
A subitánea et impro\textit{vísa} \textbf{mor}te,\hfill líbera nos Dómine.\\
Ab insídi\textit{is di}\textbf{á}boli,\hfill líbera nos Dómine.\\
Ab ira, et ódio, et omni mala \textit{volun}\textbf{tá}te,\hfill líbera nos Dómine.\\
A spíritu forni\textit{cati}\textbf{ó}nis,\hfill líbera nos Dómine.\\
A fúlgure et \textit{tempes}\textbf{tá}te,\hfill líbera nos Dómine.\\
A flagéllo \textit{terræ}\textbf{mó}tus,\hfill líbera nos Dómine.\\
A peste, fa\textit{me, et} \textbf{bel}lo,\hfill líbera nos Dómine.

\myrule

\noindent Ab imminénti\textit{bus} \textit{pér}iculis,\footnote{\rubrique{In Oratione Quadraginta Horarum tantum. At the Forty Hours' \thinspace Prayer only.}} \hfill líbera nos Dómine.
%\negativebaselineskipvspace
\myrule
\noindent A mor\textit{te per}\textbf{pé}tua,\hfill líbera nos Dómine.\\
Per mystérium sanctæ incarnati\tinyhspace\textit{ónis} \textbf{tu}æ,\hfill líbera nos Dómine.\\
Per ad\textit{véntum} \textbf{tu}um,\hfill líbera nos Dómine.\\
Per nativi\tinyhspace\textit{tátem} \textbf{tu}am,\hfill líbera nos Dómine.\\
Per baptísmum et sanctum jejú\textit{nium} \textbf{tu}um,\hfill líbera nos Dómine.\\
Per crucem et passi\tinyhspace\textit{ónem} \textbf{tu}am,\hfill líbera nos Dómine.\\
Per mortem et sepul\tinyhspace\textit{túram} \textbf{tu}am,\hfill líbera nos Dómine.\\
Per sanctam resurrecti\tinyhspace\textit{ónem} \textbf{tu}am,\hfill líbera nos Dómine.\\
Per admirábilem ascensi\tinyhspace\textit{ónem} \textbf{tu}am,\hfill líbera nos Dómine.\\
Per advéntum Spíritus San\tinyhspace\textit{cti Pa}\textbf{rá}cliti,\hfill líbera nos Dómine.\\
In di\tinyhspace\textit{e ju}\textbf{dí}cii,\hfill líbera nos Dómine.
%\end{nstabbing}

\smallscore{litaniae_sanctorum_7_peccatores_solesmes}%
%\begin{nstabbing}
%Ut nó\textit{bis} \textbf{pár}cas, te rógamus áudi nos.\\
%Ut Ecclésiam túam san\-ctam ' régere et conserváre \textit{di}\textbf{gne}ris,\\
%\hspace*{1cm} te rógamus áudi nos.
%\end{nstabbing}
\noindent
Ut no\textit{bis} \textbf{par}cas,\hfill te rogámus audi nos.\\
Ut nobis \textit{in}\textbf{dúl}geas,\hfill te rogámus audi nos.\\
Ut ad veram pæniténtiam nos perdúcere
\textit{di}\textbf{gné}ris,\hfill te rogámus audi nos.\\
Ut Ecclésiam tuam sanctam ’ régere et
conserváre\\
\emspace \textit{di}\textbf{gné}ris,\hfill te rogámus audi nos.\\
Ut Domnum Apostólicum et omnes
ecclesiásticos \\
\emspace órdines ’ in sancta  religióne conserváre \textit{di}\textbf{gné}ris,\hfill te rogámus audi nos.\\
Ut inimícos sanctæ Ecclésiæ ’ humiliáre
\textit{di}\textbf{gné}ris,\hfill te rogámus audi nos.\\
Ut régibus et princípibus christiánis ’ pacem et \\
\emspace veram concórdiam donáre \textit{di}\textbf{gné}ris,\hfill te rogámus audi nos.\\
Ut cuncto pópulo christiáno ’ pacem et
unitátem largíri \textit{di}\textbf{gné}ris,\hfill te rogámus audi nos.\\
Ut omnes errántes ad unitátem Ecclésiæ \\
\emspace
revocáre, ’ et infidéles univérsos ad Evangélii
lumen \\
\emspace perdúcere \textit{di}\textbf{gné}ris,\hfill te rogámus audi nos.\\
Ut nosmetípsos in tuo sancto servítio ’
confortáre \\
\emspace  et conserváre \textit{di}\textbf{gné}ris,\hfill te rogámus audi nos.\\
Ut mentes nostras ’ ad cæléstia desidéri\tinyhspace\textit{a}
\textbf{é}rigas,\hfill te rogámus audi nos.\\
Ut ómnibus benefactóribus nostris ’
sempitérna bona \textit{re}\textbf{trí}buas,\hfill te rogámus audi nos.\\
Ut ánimas nostras ’ fratrum, propinquórum et \\
\emspace benefactórum nostrórum ’ ab ætérna
damnatióne \textit{e}\textbf{rí}pias,\hfill te rogámus audi nos.\\
Ut fructus terræ ’ dare et conserváre
\textit{di}\textbf{gné}ris,\hfill te rogámus audi nos.\\
Ut ómnibus fidélibus defúnctis ’ réquiem
ætérnam \\
\emspace donáre \textit{di}\textbf{gné}ris,\hfill te rogámus audi nos.\\
Ut nos exaudíre \textit{di}\textbf{gné}ris,\hfill te rogámus audi nos.\\
Fi\tinyhspace\textit{li} \textbf{De}i,\hfill te rogámus audi nos.

\smallscore{litaniae_sanctorum_8_agnus_Dei_solesmes}

\smallscore{litaniae_sanctorum_9_kyrie_eleison_english_solesmes}

\textes{litaniae_san_ps_vv}

\textes{litaniae_sanctorum_orationes}
\end{document}